
\begin{definition}{7.4.26}\label{def:7-4-26}
对
$\varphi,\varphi'\in
\operatorname{Der}B(C[1])$，
定义
\begin{equation}
\delta(\varphi)
=\widehat d\circ\varphi
-(-1)^{\deg\varphi}
\varphi\circ\widehat d,
\tag{7.4.27}\label{eq:7-4-27}
\end{equation}
其中 $\widehat d$ 如定义~\ref{def:3-2-16}；
另定义
\begin{equation}
[\varphi,\varphi']
=-\frac12\Bigl(
\varphi\circ\varphi'
-(-1)^{\deg\varphi\,\deg\varphi'}
\varphi'\circ\varphi\Bigr).
\tag{7.4.28}\label{eq:7-4-28}
\end{equation}
\end{definition}

式~\eqref{eq:7-4-28}
中这个不太常见的负号，
将在命题 7.4.34、定理 7.4.41
及引理 7.4.113 中得到解释。
还须指出，这里的系数 $\frac12$
来自第~\ref{sec:3-8} 节
把 $E_k(C[1])$ 定义为
$B_k(C[1])$ 的对称群不变部分。
在其他一些文献中，
$E_k(C[1])$ 则定义为
$B_k(C[1])$ 对由形如
\[
\cdots x\otimes y\cdots
-(-1)^{\deg' x\,\deg' y}
(\cdots y\otimes x\cdots)
\]
的元素生成的理想之商。
在此约定下，
$x\otimes y\pm y\otimes x$
与 $[x\otimes y]$
（即 $x\otimes y$ 的等价类）
之间相差一个 $\frac12$ 的系数。
另回顾，本书取
$e^x=\sum_k x^{\otimes k}$，
而不取
$e^x=\sum_k \frac1{k!}x^{\otimes k}$。

借助
\eqref{eq:7-4-25-1}
将
\[
\varphi^i=(\varphi_k^i)_{k\in\mathbb Z_{\ge0}}
\in\prod_{k=0}^{\infty}
\operatorname{Hom}^{\mathrm{filt}}
(B_k(C[1]),C[1])
\cong\operatorname{Der}B(C[1])
\]
等同起来，
则 $\delta(\varphi)
=((\delta\varphi)_k)$ 满足
\begin{equation}
\begin{aligned}
(\delta\varphi)_k(\mathbf x)
={}&\sum_c
(-1)^{\deg\varphi\,\deg'\mathbf x_c^{(3;1)}}
\mathfrak m\bigl(
\mathbf x_c^{(3;1)}\otimes
\varphi(\mathbf x_c^{(3;2)})\otimes
\mathbf x_c^{(3;3)}\bigr)\\
&-\sum_c
(-1)^{\deg\varphi+
\deg'\mathbf x_c^{(3;1)}}
\varphi\bigl(
\mathbf x_c^{(3;1)}\otimes
\mathfrak m(\mathbf x_c^{(3;2)})\otimes
\mathbf x_c^{(3;3)}\bigr).
\end{aligned}
\tag{7.4.27$'$}\label{eq:7-4-27-prime}
\end{equation}
并且
\begin{equation}
\begin{aligned}
([\varphi^1,\varphi^2])_k(\mathbf x)
={}&\frac12\sum_c
(-1)^{\deg\varphi^2\,\deg'\mathbf x_c^{(3;1)}+1}
\varphi^1\bigl(
\mathbf x_c^{(3;1)}\otimes
\varphi^2(\mathbf x_c^{(3;2)})\\
&\hspace{5em}\otimes
\mathbf x_c^{(3;3)}\bigr)\\
&+\frac12\sum_c
(-1)^{\deg\varphi^1
(\deg'\mathbf x_c^{(3;1)}
+\deg\varphi^2)}
\varphi^2\bigl(
\mathbf x_c^{(3;1)}\otimes
\varphi^1(\mathbf x_c^{(3;2)})\\
&\hspace{5em}\otimes
\mathbf x_c^{(3;3)}\bigr).
\end{aligned}
\tag{7.4.28$'$}\label{eq:7-4-28-prime}
\end{equation}
这里若
$\mathbf x=x_1\otimes\cdots\otimes x_\ell$，
则约定
$\deg'\mathbf x
=\sum_i\deg' x_i
=\deg\mathbf x-\ell$。

\begin{lemma}{7.4.29}\label{lem:7-4-29}
式~\eqref{eq:7-4-27}
与~\eqref{eq:7-4-28}
的右端均属于
$\operatorname{Der}B(C[1])$。
此外，
$\bigl(\operatorname{Der}B(C[1]),
\delta,[\ ,\ ]\bigr)$
是微分分次李代数。
\end{lemma}

证明直接且标准，故从略。
再把后文第 A3 节例 A3.4
所给构造取滤过版本，
便使 $\operatorname{Der}B(C[1])$
成为滤过 $L_\infty$ 代数。

现在研究一个严格滤过
$L_\infty$ 同态
\begin{equation}
\widehat{\mathfrak q}^{\circ}:
C^{\mathrm{amb}}[1]
\longrightarrow\operatorname{Der}B(C[1]).
\tag{7.4.30}\label{eq:7-4-30}
\end{equation}
这里上标 $\circ$ 表示“反序”\liftIndex{I288}，
见式 (7.4.39.1)。
从现在起，
我们假定所论的一切滤过
$L_\infty$ 同态
均严格、滤过且有隙。

按严格滤过 $L_\infty$ 同态的定义，
$\widehat{\mathfrak q}^{\circ}$
是一列滤过同态：
\begin{equation}
\widehat{\mathfrak q}^{\circ}_\ell:
E_\ell(C^{\mathrm{amb}}[2])
\longrightarrow
(\operatorname{Der}B(C[1]))[1],
\qquad \ell=1,2,\ldots,
\tag{7.4.31}\label{eq:7-4-31}
\end{equation}
它们的次数为 $0$，
并满足一个适当的关系，
该关系稍后将在命题 7.4.34 讨论。
注意，因为
$\widehat{\mathfrak q}^{\circ}$
是严格的，
只有 $\ell\ge1$ 的算子
$\mathfrak q^{\circ}_\ell$，
而没有 $\ell=0$ 的分量。

经
$CH(C,C)\cong\operatorname{Der}B(C[1])$
这一识别\translatornote{源书此处把
$\operatorname{Der}B(C[1])$
写作多一枚右括号的
$\operatorname{Der}B(C[1]))$；
这里依前文已定义的记号排印，
见疑误 A-0563。}，
一列式~\eqref{eq:7-4-31}
的滤过同态
等价于一族滤过同态
\begin{equation}
\mathfrak q^{\circ}_{\ell,k}:
E_\ell(C^{\mathrm{amb}}[2])
\otimes B_k(C[1])
\longrightarrow C[1],
\quad \ell=1,2,\ldots,\quad k=0,1,\ldots,
\tag{7.4.32}\label{eq:7-4-32}
\end{equation}
后者的次数为 $1$。
另置
\begin{equation}
\mathfrak q^{\circ}_{0,k}
=\mathfrak m_k:
B_k(C[1])\longrightarrow C[1].
\tag{7.4.33}\label{eq:7-4-33}
\end{equation}
